\documentclass[11pt]{article}
\usepackage[margin=1in]{geometry}
\usepackage{amsmath,amssymb}

\title{Core Math Behind \texttt{vicon\_goto\_mppi\_controller.py}}
\author{}
\date{\today}

\begin{document}
\maketitle

\section{Purpose}

This note explains the core control math implemented in
\texttt{scripts/vicon\_goto\_mppi\_controller.py}. The goal of that script is
to drive a differential-drive JetBot from its current Vicon pose to a planar
target $(x_d, y_d)$ by repeatedly solving a short finite-horizon control
problem with Model Predictive Path Integral (MPPI) control.

The main code blocks discussed here are:
\begin{itemize}
\item heuristic command seed: lines 236--264,
\item rollout dynamics model: lines 266--283,
\item rollout cost: lines 285--326,
\item heuristic rollout construction: lines 328--343,
\item MPPI update: lines 345--419,
\item outer control loop: lines 421--489.
\end{itemize}

\section{State, Target, and Heading}

In the outer loop (lines 443--456), the script reads the latest Vicon pose and
computes:
\[
\Delta x = x_d - x,
\qquad
\Delta y = y_d - y,
\qquad
r = \sqrt{(\Delta x)^2 + (\Delta y)^2}.
\]

The target heading is
\[
\theta_d = \operatorname{atan2}(\Delta y, \Delta x).
\]

The robot heading is obtained from the Vicon orientation by rotating the body
$Y$ axis into the world frame, projecting onto the floor plane, and applying a
user-entered angle correction. In compact form,
\[
\theta = \operatorname{wrapToPi}\!\left(\operatorname{atan2}(f_y, f_x) + \delta\right),
\]
where $(f_x,f_y)$ is the horizontal forward vector and $\delta$ is the angle
correction. The heading error is then
\[
e_\theta = \operatorname{wrapToPi}(\theta_d - \theta).
\]

If $r \le \varepsilon$, the controller stops immediately (lines 458--469).

\section{Heuristic Command Used as the MPPI Seed}

Before sampling random control perturbations, the script builds a simple
goal-seeking differential-drive command in \texttt{\_heuristic\_command()}
(lines 236--264). This heuristic is not the final controller by itself; it is
the warm-start policy that biases MPPI toward sensible commands.

\subsection{Distance-Based Forward Speed}

Let:
\begin{itemize}
\item $v_{\max} =$ \texttt{self.\_speed},
\item $v_{\min} =$ \texttt{self.\_min\_forward\_speed},
\item $r_s =$ \texttt{self.\_slow\_radius}.
\end{itemize}

Then the base forward command is
\[
v_{\mathrm{base}} =
\begin{cases}
v_{\max}, & r \ge r_s, \\
v_{\min} + (v_{\max} - v_{\min}) \dfrac{r}{r_s}, & r < r_s.
\end{cases}
\]

This matches lines 248--253.

\subsection{Heading Alignment and Steering}

The same block computes
\[
a = \max(0, \cos e_\theta),
\qquad
v = v_{\mathrm{base}} a^2,
\]
so the robot reduces forward motion when it is poorly aligned with the target
(lines 255--258).

Steering is proportional to heading error and then clamped:
\[
\omega = \operatorname{clamp}\!\left(k_\theta e_\theta,\; -u_{\mathrm{turn}},\; u_{\mathrm{turn}}\right),
\]
where $k_\theta =$ \texttt{self.\_heading\_gain} and
$u_{\mathrm{turn}} =$ \texttt{self.\_turn\_speed} (lines 259--263).

The heuristic left and right wheel commands are
\[
u_L = \operatorname{clamp}(v - \omega),
\qquad
u_R = \operatorname{clamp}(v + \omega),
\]
as returned on line 264. Here the one-argument \texttt{clamp()} means clamp to
the motor command range, effectively $[-1,1]$.

\section{Internal Rollout Dynamics Model}

The MPPI controller needs a fast predictive model to evaluate candidate command
sequences. The model in \texttt{\_simulate\_step()} (lines 266--283) uses a
simple differential-drive abstraction.

For wheel commands $(u_L, u_R)$:
\[
u_f = \frac{u_L + u_R}{2},
\qquad
u_\omega = \frac{u_R - u_L}{2}.
\]

These are interpreted as forward and turn commands. They are converted into
physical rates by calibration constants
\texttt{self.\_model\_linear\_speed} and
\texttt{self.\_model\_angular\_speed}:
\[
v = c_v u_f,
\qquad
\omega = c_\omega u_\omega.
\]

Over one controller period $\Delta t =$ \texttt{self.\_period}, the predicted
state update is
\begin{align*}
x_{k+1} &= x_k + v \cos(\theta_k)\,\Delta t, \\
y_{k+1} &= y_k + v \sin(\theta_k)\,\Delta t, \\
\theta_{k+1} &= \operatorname{wrapToPi}(\theta_k + \omega \Delta t).
\end{align*}

This is exactly what lines 274--282 implement.

\section{Rollout Cost Function}

The cost of a candidate control sequence is computed in
\texttt{\_sequence\_cost()} (lines 285--326). Suppose a rollout consists of
\[
\mathbf{u}_{0:H-1} = \left\{(u_{L,0},u_{R,0}), \dots, (u_{L,H-1},u_{R,H-1})\right\}.
\]

At each rollout step $k$, the script simulates the next state and evaluates:
\begin{itemize}
\item distance-to-goal penalty,
\item heading-error penalty,
\item control-effort penalty,
\item control-slew penalty relative to the previous command,
\item reverse-motion penalty.
\end{itemize}

\subsection{Stage Cost}

Let
\[
r_k = \sqrt{(x_d - x_k)^2 + (y_d - y_k)^2},
\qquad
e_{\theta,k} = \operatorname{wrapToPi}\!\left(\operatorname{atan2}(y_d-y_k, x_d-x_k) - \theta_k\right).
\]

The code scales distance by
\[
\bar{r}_k = \frac{r_k}{r_{\mathrm{scale}}},
\qquad
r_{\mathrm{scale}} = \texttt{self.\_distance\_scale},
\]
with \texttt{self.\_distance\_scale} initialized from the slow radius.

Let the previous command be $(u_{L,k-1},u_{R,k-1})$. Then the stage cost used
by lines 308--315 can be written as
\begin{align*}
\ell_k
&=
w_r \bar{r}_k^2
+ w_\theta e_{\theta,k}^2
+ w_u \left(u_{L,k}^2 + u_{R,k}^2\right) \\
&\quad
+ w_{\Delta u} \left((u_{L,k} - u_{L,k-1})^2 + (u_{R,k} - u_{R,k-1})^2\right)
+ w_{\mathrm{rev}} \max(-u_{f,k}, 0)^2,
\end{align*}
where
\[
u_{f,k} = \frac{u_{L,k} + u_{R,k}}{2}.
\]

The weights above correspond to:
\begin{itemize}
\item \texttt{self.\_distance\_cost\_weight},
\item \texttt{self.\_heading\_cost\_weight},
\item \texttt{self.\_control\_effort\_weight},
\item \texttt{self.\_control\_slew\_weight},
\item \texttt{self.\_reverse\_cost\_weight}.
\end{itemize}

\subsection{Terminal Cost}

After the final simulated step, lines 319--325 add a terminal penalty:
\[
\phi =
w_r^{(T)} \bar{r}_H^2
+ w_\theta^{(T)} e_{\theta,H}^2,
\]
using
\texttt{self.\_terminal\_distance\_cost\_weight} and
\texttt{self.\_terminal\_heading\_cost\_weight}.

The full rollout objective is therefore
\[
J(\mathbf{u}_{0:H-1}) = \sum_{k=0}^{H-1} \ell_k + \phi.
\]

\section{How the Heuristic Becomes a Nominal Sequence}

Lines 328--343 build a horizon-length heuristic sequence by repeatedly applying
the heuristic command and rolling the internal model forward:
\[
\mathbf{u}^{\mathrm{heur}}_{0:H-1}
=
\left\{
\pi_{\mathrm{heur}}(x_0,y_0,\theta_0),
\pi_{\mathrm{heur}}(x_1,y_1,\theta_1),
\dots
\right\}.
\]

Then lines 354--361 combine:
\begin{itemize}
\item the previous cycle's shifted nominal sequence, and
\item the fresh heuristic sequence.
\end{itemize}

If $\alpha =$ \texttt{self.\_heuristic\_blend}, the nominal sequence at step
$k$ is:
\[
\mathbf{u}^{\mathrm{nom}}_k
=
(1-\alpha)\,\mathbf{u}^{\mathrm{shift}}_k
+ \alpha\,\mathbf{u}^{\mathrm{heur}}_k.
\]

This warm start is important: it makes the sampled rollouts less random and
gives MPPI a plausible local baseline.

\section{MPPI Sampling and Softmin Update}

The core MPPI logic lives in \texttt{\_mppi\_command()} (lines 345--419).

\subsection{Sampling Candidate Sequences}

For each sample index $i \in \{1,\dots,N\}$ and horizon step $k$, the script
draws Gaussian noise
\[
\epsilon^i_k =
\begin{bmatrix}
\epsilon^i_{L,k} \\
\epsilon^i_{R,k}
\end{bmatrix},
\qquad
\epsilon^i_{L,k}, \epsilon^i_{R,k}
\sim \mathcal{N}(0,\sigma^2),
\]
where $\sigma =$ \texttt{self.\_noise\_std} (lines 370--379).

The sampled command is
\[
\mathbf{u}^i_k
=
\operatorname{clamp}\!\left(\mathbf{u}^{\mathrm{nom}}_k + \epsilon^i_k\right).
\]

Each full candidate sequence is then scored by the rollout cost
\[
J_i = J(\mathbf{u}^i_{0:H-1}),
\]
as in lines 381--389.

\subsection{Exponential Weighting}

After evaluating all candidates, the controller computes the minimum cost
\[
J_{\min} = \min_i J_i
\]
and then assigns each rollout a softmin weight (lines 393--395):
\[
w_i = \exp\!\left(-\frac{J_i - J_{\min}}{\lambda}\right),
\]
where $\lambda =$ \texttt{self.\_temperature}.

Subtracting $J_{\min}$ does not change the relative weighting, but it improves
numerical stability by keeping exponents closer to zero.

\subsection{Weighted Noise Update}

For each horizon step $k$, the script computes the weighted mean noise:
\[
\bar{\epsilon}_k =
\frac{\sum_{i=1}^N w_i \epsilon^i_k}{\sum_{i=1}^N w_i}.
\]

The nominal sequence is updated by
\[
\mathbf{u}^{\mathrm{new}}_k
=
\operatorname{clamp}\!\left(\mathbf{u}^{\mathrm{nom}}_k + \bar{\epsilon}_k\right),
\]
which is what lines 400--416 implement.

Finally, line 419 returns only the first control:
\[
\mathbf{u}_{\mathrm{apply}} = \mathbf{u}^{\mathrm{new}}_0.
\]

This is the standard receding-horizon idea: optimize a short sequence, apply
the first command, shift the horizon, and repeat on the next cycle.

\section{Outer Controller Loop}

The outer loop in \texttt{run()} (lines 421--489) connects the math to the
live robot:
\begin{enumerate}
\item get the active target and current tuning (lines 423--424),
\item wait if no target is active (lines 426--428),
\item read the latest Vicon pose and verify visibility (lines 430--441),
\item compute current planar state and heading error (lines 443--456),
\item stop if the robot is already within $\varepsilon$ (lines 458--469),
\item read the currently commanded wheel speeds (line 472),
\item call \texttt{\_mppi\_command(...)} to get the next wheel command
      (lines 473--480),
\item send the command to the JetBot and update status (lines 482--488).
\end{enumerate}

The current wheel command is passed into the cost function so the
control-slew term can discourage abrupt changes.

\section{Compact Summary}

The controller can be summarized as:
\begin{align*}
\text{measure current state} &\to (x,y,\theta), \\
\text{build heuristic seed} &\to \mathbf{u}^{\mathrm{heur}}_{0:H-1}, \\
\text{blend with shifted previous plan} &\to \mathbf{u}^{\mathrm{nom}}_{0:H-1}, \\
\text{sample noisy sequences} &\to \mathbf{u}^i_{0:H-1}, \\
\text{score each sequence} &\to J_i, \\
\text{compute softmin weights} &\to w_i, \\
\text{update the nominal plan} &\to \mathbf{u}^{\mathrm{new}}_{0:H-1}, \\
\text{apply the first command} &\to (u_L, u_R).
\end{align*}

So the script is not using an analytic optimal-control solution. Instead, it
uses a simple dynamics model plus many short random control perturbations, then
keeps the perturbation directions that tend to reduce the finite-horizon cost.

\section{Parameter Anchors in the File}

The MPPI-related tunable parameters are declared near lines 526--550:
\begin{itemize}
\item speed and turning limits: lines 526--533,
\item horizon and sample count: lines 537--540,
\item model calibration: lines 541--542,
\item heuristic blending: line 543,
\item distance, heading, effort, slew, reverse, and terminal weights:
      lines 544--550.
\end{itemize}

Those parser defaults determine the numerical behavior of the equations above.

\end{document}
